\documentclass[aps,prl,reprint,superscriptaddress,nofootinbib]{revtex4-2}

\usepackage{amsmath,amssymb,mathtools,bm}
\usepackage{microtype}
\usepackage[hidelinks]{hyperref}

\newcommand{\Tr}{\operatorname{Tr}}
\newcommand{\diag}{\operatorname{diag}}
\newcommand{\spec}{\operatorname{spec}}
\newcommand{\e}{\mathrm{e}}
\newcommand{\id}{\mathbb{I}}
\newcommand{\C}{\mathbb{C}}
\newcommand{\Z}{\mathbb{Z}}
\newcommand{\ket}[1]{\lvert #1\rangle}
\newcommand{\bra}[1]{\langle #1\rvert}
\newcommand{\PhiD}{\lvert\Phi_d\rangle}
\newcommand{\Phithree}{\lvert\Phi_3\rangle}

\begin{document}

\title{Optimality of CGLMP Measurements for Maximally Entangled Qutrits}

\author{Gurvir Singh}
\email{gurvirsingh@hri.res.in}
\affiliation{Harish-Chandra Research Institute, A CI of Homi Bhabha National Institute, Chhatnag Road, Jhunsi, Prayagraj 211019, India}

\date{\today}

\begin{abstract}
For the maximally entangled state, the standard Fourier-phase measurements have long been conjectured to give the largest violation of the Collins-Gisin-Linden-Massar-Popescu inequality. We prove this conjecture exactly for qutrits, with no restriction on the local projective measurements. The proof reveals a geometric structure behind the problem. The four relative measurement bases form a closed unitary path whose holonomy is the cyclic shift. In dimension three, the CGLMP cost of each step is completely encoded by a single commutator in $SU(3)$. The shift imposes a fixed total angular displacement, and the cost is minimized when this displacement is shared equally among the four steps. The standard CGLMP measurements realize precisely this balanced factorization. Consequently,
\[
I_3(\Phi_3)\leq \frac{4}{3}+\frac{8}{3\sqrt3}=2.872934051\ldots,
\]
with equality for the usual CGLMP settings. We also extend the holonomy analysis to arbitrary dimension, establishing simultaneous angular-budget inequalities and equal sharing along the Fourier path, while identifying the remaining obstruction to a global geometric proof.
\end{abstract}

\maketitle

\paragraph{Introduction.}
Bell inequalities with more than two outcomes probe structures that are absent in the Clauser-Horne-Shimony-Holt scenario. The Collins-Gisin-Linden-Massar-Popescu (CGLMP) family is the canonical example \cite{CGLMP}; its inequalities are facets of the local correlation polytope \cite{Masanes}. The construction is built from cyclic differences of outcomes. For a deterministic local model, three differences around the four-setting cycle may be chosen independently, while the fourth is constrained by closure of the cycle \cite{CGLMP}. This simple frustration mechanism is one reason the CGLMP family has remained a useful laboratory for high-dimensional nonlocality.

The quantum optimization already displays a feature with no qubit analogue. For the maximally entangled state
\begin{equation}
\PhiD=\frac{1}{\sqrt d}\sum_{j=0}^{d-1}\ket{j,j},
\label{eq:mes}
\end{equation}
the standard Fourier-phase measurements introduced in the early studies of high-dimensional Bell tests \cite{Durt2001,CGLMP} give, for $d=3$, the value $I_3=4/3+8/(3\sqrt3)\simeq2.8729$. Ac\'{i}n, Durt, Gisin, and Latorre then showed that the same measurement family produces the larger value $1+\sqrt{11/3}\simeq2.9149$ when the state is allowed to become nonmaximally entangled \cite{Acin2002}. The distinction is essential: the unrestricted Tsirelson optimization and the measurement optimization at fixed maximally entangled state are different problems. Later work reinforced the role of nonmaximally entangled states in optimal Bell tests \cite{AcinGillGisin}, and exact sum-of-squares certificates now determine the unrestricted CGLMP quantum maxima for $d=3,4$ \cite{IoannouRosset}.

The fixed-state problem has remained much more resistant. Zohren and Gill explicitly identified the Fourier-phase settings as the conjectured optimal measurements for $\PhiD$ and found strong numerical support over a wide range of dimensions \cite{ZG}. Semidefinite hierarchies provide another route to upper bounds on Bell violations \cite{NPA}. In particular, the hierarchy for correlations generated by maximally entangled states introduced by Lang, V\'{e}rtesi, and Navascu\'{e}s reproduces the qutrit candidate to roughly $10^{-10}$ numerical accuracy \cite{LVN}. A recent all-dimensional proof establishes the optimal projective value and classifies its maximizers by different analytic methods \cite{TurtureanBao}. The question pursued here is whether the qutrit result admits a direct geometric explanation and how far that explanation extends in dimension.

There is also a complementary line of work in which the Bell functional, rather than the state, is redesigned. Salavrakos \textit{et al.} constructed Bell inequalities whose exact quantum maximum is attained by maximally entangled qudits together with the same CGLMP measurement family \cite{SATWAP}; those correlations were later used for self-testing maximally entangled states of arbitrary local dimension with only two measurements per party \cite{Sarkar}. Thus the Fourier-phase measurements have repeatedly appeared as a distinguished high-dimensional measurement geometry. The original CGLMP facet, however, calls for an optimization at fixed maximally entangled state rather than a tailored Bell functional.

Here we give an independent geometric proof for qutrits. The argument does not solve stationarity equations for a parameterization of $U(3)$. Instead, fixing $\Phithree$ exposes a change of variables in which the four relative measurement bases form a unitary loop with prescribed holonomy. In dimension three, the Bell cost of each edge is encoded by one $SU(3)$ commutator. The holonomy imposes a fixed angular budget, and a convexity argument shows that the minimum is attained by sharing this budget equally among all four edges. This viewpoint also yields two consequences beyond the value of the optimum: every optimal strategy has the same edgewise modular-displacement profile, and the standard settings are exactly optimal for white-noise robustness of $\Phithree$ within qutrit projective measurements.

Our main result is the following.

\textit{Theorem.} Let Alice and Bob share $\Phithree$ and let each party perform two arbitrary complete rank-one qutrit projective measurements. Then the qutrit CGLMP value satisfies
\begin{equation}
I_3\leq \frac{4}{3}+\frac{8}{3\sqrt3}.
\label{eq:mainI}
\end{equation}
Equality is attained by the standard Fourier-phase CGLMP measurements.

Thus the conjectured measurements are globally optimal for the fixed maximally entangled qutrit state. No Fourier, interferometric, symmetry, or reality assumption is imposed on the competing measurements.

\paragraph{From four measurement bases to one unitary loop.}
It is convenient to use the modular form of CGLMP. For outcomes in $\Z_d$, write $[r]_d\in\{0,\ldots,d-1\}$ for the residue of $r$ and define
\begin{align}
B_d={}&\mathbb E[A_0-B_0]_d+\mathbb E[B_0-A_1]_d \notag\\
&+\mathbb E[A_1-B_1]_d+\mathbb E[B_1-A_0-1]_d .
\label{eq:Bd}
\end{align}
Local correlations obey $B_d\geq d-1$. For $d=3$, the usual normalization is $I_3=4-B_3$.

Let the columns of $\mathsf A_x$ and $\mathsf B_y$ be the local measurement bases. On the maximally entangled state, every joint probability is an overlap of two basis vectors. This allows the four terms in Eq.~\eqref{eq:Bd} to be expressed through a single directed cost
\begin{equation}
c_d(U)=\frac{1}{d}\sum_{r,s=0}^{d-1}[r-s]_d\,|U_{rs}|^2.
\label{eq:cost}
\end{equation}
If $S\ket j=\ket{j-1\;\mathrm{mod}\;d}$ is the cyclic lowering shift, the measurement optimization is exactly equivalent to
\begin{equation}
B_d=\sum_{j=1}^{4}c_d(U_j),\qquad U_1U_2U_3U_4=S.
\label{eq:factor}
\end{equation}
The explicit map between the four measurement bases and the four unitaries, together with the converse construction, is given in the Supplemental material.

Equation~\eqref{eq:factor} is the first key point. The four measurement choices do not produce four independent overlap matrices. They form a path whose accumulated mismatch is the cyclic shift. There is a direct conceptual parallel with the original local-variable derivation of CGLMP. Classically, the four modular differences around the cycle must close additively, which forces one relation to be frustrated \cite{CGLMP}. For the fixed maximally entangled quantum problem, the corresponding closure law is multiplicative and noncommutative: the four relative bases multiply to $S$ rather than to the identity. In this sense the shift in Eq.~\eqref{eq:factor} is a unitary holonomy carrying the same cyclic frustration into measurement space. The standard CGLMP settings correspond to four identical quarter-steps of this holonomy.

\paragraph{Proof idea.}
The proof has four steps. First, the qutrit displacement statistics hidden in each relative unitary are compressed into one multiplicative commutator. Second, the four commutators inherit the holonomy of the unitary loop and multiply to a fixed central element of $SU(3)$. Third, that central element imposes a minimum total angular displacement. Finally, a sharp one-matrix bound shows that the cheapest way to realize this displacement is to divide it equally among the four steps. The technical derivations are collected in the Supplemental material; here we keep the geometric content visible.

Set $d=3$, $\omega=\e^{2\pi i/3}$, and $Z=\diag(1,\omega,\omega^2)$. To every relative unitary $U$ associate
\begin{equation}
K(U)=ZUZ^\dagger U^\dagger\in SU(3).
\label{eq:K}
\end{equation}
The three cyclic displacement probabilities carried by $|U_{rs}|^2$ have only two independent real parameters. Those two parameters are exactly the real and imaginary parts of $\Tr K(U)$. Consequently, the qutrit cost $c_3(U)$ is an exact conjugation-invariant function $\ell(K(U))$, not an approximation or a bound. This collapse is the first genuinely qutrit-specific ingredient.

The map $K$ obeys a cocycle identity. Transporting the four commutators along the successive partial products of the path in Eq.~\eqref{eq:factor} produces matrices $L_j\in SU(3)$ satisfying
\begin{equation}
L_1L_2L_3L_4=\omega^2\id,
\qquad B_3=\sum_{j=1}^{4}\ell(L_j).
\label{eq:central}
\end{equation}
The Bell problem has therefore become a cost minimization over four $SU(3)$ elements with a fixed nontrivial central product. The original measurement coordinates have disappeared.

To quantify the obstruction carried by this product, define $m(K)$ as the smallest possible largest eigenangle among traceless Hermitian logarithms of $K$. In the $SU(3)$ Weyl alcove this is simply the largest canonical eigenangle. Thompson's product-of-exponentials theorem \cite{Thompson,Antezana} makes $m$ subadditive. Since $m(\omega^2\id)=4\pi/3$, Eq.~\eqref{eq:central} forces the angular budget
\begin{equation}
\sum_{j=1}^{4}m(L_j)\geq\frac{4\pi}{3}.
\label{eq:angbudget}
\end{equation}

The remaining question is how much Bell cost is required to carry a given amount of this angular budget. Minimizing $\ell(K)$ at fixed $m(K)$ gives a sharp scalar envelope $g(m)$. On the branch containing the optimum, $g$ is increasing and strictly convex. A second branch appears when a single factor carries angular displacement at least $2\pi/3$, but it is strictly more expensive than the balanced configuration. The Supplemental material proves these statements by a one-variable minimization in the Weyl alcove followed by an elementary scalar estimate.

The geometry is therefore simple. The product constraint demands a total displacement of at least $4\pi/3$, while convexity penalizes concentrating that displacement in one or two factors. The least expensive possibility is to share it equally, so each factor carries $m=\pi/3$. This gives
\begin{equation}
B_3\geq\frac{8}{3}\left(1-\frac{1}{\sqrt3}\right),
\label{eq:Bbound}
\end{equation}
which is equivalent to Eq.~\eqref{eq:mainI}.

\paragraph{What optimality forces.}
The equality conditions contain more information than the optimal value. Strict convexity at the balanced point, together with the strict exclusion of the second branch, implies that every global optimizer must satisfy
\begin{equation}
m(L_1)=m(L_2)=m(L_3)=m(L_4)=\frac{\pi}{3}.
\label{eq:equalm}
\end{equation}
Moreover, equality in the one-matrix envelope fixes the canonical spectrum of every transported commutator, and hence also of every $K(U_j)$, to
\begin{equation}
\spec K(U_j)=\{\e^{i\pi/3},\e^{-i\pi/6},\e^{-i\pi/6}\}.
\label{eq:rigidspec}
\end{equation}
This is a partial rigidity statement for arbitrary optimal projective measurements. It does not by itself prove uniqueness of the four bases up to all CGLMP symmetries, but it shows that any optimizer must realize the same commutator spectrum on every edge of the loop.

There is an equivalent probabilistic statement. If $p_r(U_j)$ denotes the total probability weight on the $r$th cyclic diagonal of $|U_j|^2$, then the trace relation fixing $K(U_j)$ gives, for every edge,
\begin{equation}
(p_0,p_1,p_2)=\left(\frac{4+2\sqrt3}{9},\frac{4-2\sqrt3}{9},\frac19\right).
\label{eq:edgeprofile}
\end{equation}
Thus optimality is not achieved by trading a very good setting pair against a worse one. Each of the four directed CGLMP terms carries exactly the same cost. The balanced angular picture therefore has a directly observable counterpart in the modular outcome statistics.

\paragraph{Saturation by the standard phases.}
The lower bound is global because the reduction above allowed arbitrary qutrit projective bases. It remains to identify a physical strategy that reaches it. Let $F$ be the qutrit Fourier matrix and write $R(t)=F\diag(1,\omega^t,\omega^{2t})F^\dagger$. Since $R(1)=S$, four equal quarter-steps close the required loop. The usual CGLMP phases can be chosen so that
\begin{equation}
U_1=U_2=U_3=U_4=R(1/4).
\label{eq:quarter}
\end{equation}
For this quarter-step, the associated commutator has canonical eigenangles $\pi/3,-\pi/6,-\pi/6$. Thus every factor carries exactly one quarter of the required angular budget and saturates the one-matrix cost bound. The standard CGLMP measurements therefore attain Eq.~\eqref{eq:Bbound} exactly.

This gives a direct interpretation of the fractional Fourier phases. They are not merely a successful ansatz. For the maximally entangled qutrit state they implement the equal-step factorization of the cyclic holonomy enforced by the Bell expression.

A useful immediate consequence concerns white noise. For
\begin{equation}
\rho_v=v\ket{\Phi_3}\!\bra{\Phi_3}+(1-v)\frac{\id}{9},
\end{equation}
the maximally mixed part contributes zero to $I_3$, so for any fixed projective measurements $I_3(\rho_v)=vI_3(\Phi_3)$. The theorem therefore also gives the exact best critical visibility
\begin{equation}
v_{\rm crit}=\frac{2}{I_3^{\max}(\Phi_3)}=3\sqrt3-\frac92\simeq0.696152.
\label{eq:vcrit}
\end{equation}
This value was obtained for the standard settings in the early qutrit analysis \cite{Acin2002}; the present result shows that no other qutrit projective measurements can improve it for the fixed maximally entangled state.

\paragraph{Higher dimensions: coupled angular budgets.}
The exact loop formulation in Eq.~\eqref{eq:factor} is not confined to qutrits. What is special to $d=3$ is the reduction of the entire cost to the spectrum of one $SU(3)$ commutator. To see the higher-dimensional structure, let $\omega_d=\e^{2\pi i/d}$, $Z_d=\diag(1,\omega_d,\ldots,\omega_d^{d-1})$, and introduce the harmonic commutators and their normalized traces,
\begin{equation}
K_k(U)=Z_d^kUZ_d^{-k}U^\dagger,\qquad
 t_k(U)=\frac{1}{d}\Tr K_k(U).
\label{eq:Kkmain}
\end{equation}
The cyclic displacement probabilities are the inverse discrete Fourier transform of the $t_k$. In particular, the exact one-edge cost is
\begin{equation}
c_d(U)=\frac{d-1}{2}-\sum_{k=1}^{d-1}
\frac{t_k(U)}{1-\omega_d^{-k}}.
\label{eq:costharmmain}
\end{equation}
Conjugate harmonics combine to make this expression real. For qutrits $t_2=\overline{t_1}$; from $d=4$ onward the cost probes genuinely distinct harmonics. They are nevertheless constrained by the \emph{same} unitary $U$, rather than being independent optimization variables.

The cocycle identity for each $K_k$ transfers the loop closure $U_1U_2U_3U_4=S$ to a central product. With $P_j=U_1\cdots U_j$, the transported commutators $L_{j,k}=P_{j-1}K_k(U_j)P_{j-1}^\dagger$ satisfy $\prod_jL_{j,k}=\omega_d^{-k}\id$. Extend the angular gauge $m$ to $SU(d)$ by minimizing the largest eigenvalue of a traceless Hermitian logarithm. Subadditivity then yields a family of necessary inequalities for \emph{every} projective measurement strategy:
\begin{equation}
\sum_{j=1}^{4}m\bigl(K_k(U_j)\bigr)
\geq 2\pi\left(1-\frac{k}{d}\right),\quad 1\leq k<d.
\label{eq:budgetdmain}
\end{equation}
The qutrit budget, Eq.~\eqref{eq:angbudget}, is the $k=1$, $d=3$ member. For $d=4$ the three right-hand sides are $3\pi/2$, $\pi$, and $\pi/2$. Thus higher-dimensional CGLMP carries a hierarchy of simultaneous angular obstructions, rather than a single scalar displacement.

\paragraph{Equal sharing survives on the Fourier path.}
Let $R_d(t)=F_d\diag(1,\omega_d^t,\ldots,\omega_d^{(d-1)t})F_d^\dagger$, where $F_d$ is the discrete Fourier matrix, so $R_d(1)=S$. Direct diagonalization gives the eigenvalues of $K_k(R_d(t))$ as $\e^{2\pi it(1-k/d)}$ with multiplicity $k$ and $\e^{-2\pi itk/d}$ with multiplicity $d-k$. Consequently, for $0\leq t\leq1$,
\begin{equation}
m\bigl(K_k(R_d(t))\bigr)=
2\pi t\left(1-\frac{k}{d}\right).
\label{eq:mFouriermain}
\end{equation}
Hence the four quarter-steps $R_d(1/4)$ simultaneously saturate \emph{all} the budgets in Eq.~\eqref{eq:budgetdmain}. Indeed, any positive subdivision $U_j=R_d(t_j)$ with $\sum_jt_j=1$ saturates them. Saturation alone therefore does not select the equal subdivision.

The cost does select it. The Supplemental material evaluates $f_d(t)=c_d(R_d(t))$ as an explicit weighted trigonometric sum and proves a sine-averaging inequality that applies even where $f_d$ ceases to be convex. For every $d\geq2$ and every $t_j\geq0$ with $\sum_jt_j=1$, one obtains the sharp \emph{Fourier-path equal-sharing inequality}
\begin{equation}
\sum_{j=1}^{4}c_d(R_d(t_j))\geq
4c_d\bigl(R_d(1/4)\bigr).
\label{eq:fouriersharemain}
\end{equation}
Thus the familiar fractional phases are the least-cost \emph{positive Fourier subdivision} of the cyclic holonomy in every dimension, not just at $d=3$. Their Bell value is the standard candidate
\begin{equation}
I_d^{\rm F}=\frac{4}{d(d-1)}\sum_{r=1}^{d-1}
(d-r)\sec\!\left(\frac{\pi r}{2d}\right).
\label{eq:IdFmain}
\end{equation}

The restriction to the Fourier path is essential. Already for $d=4$, the rotations $U=\e^{i(\pi/4)X_{01}}$ and $V=\e^{i(\pi/6)X_{02}}$, with $X_{ab}=\ket a\bra b+\ket b\bra a$, have the same $K_1$ eigenvalues $\{\e^{\pm i\pi/3},1,1\}$ but different costs: $c_4(U)=1/2$ and $c_4(V)=1/4$. Their second harmonics distinguish them. Hence no bound based only on the $K_1$ spectrum can reproduce the exact cost for arbitrary ququart measurements. A global geometric proof would require a joint cost inequality for the harmonics subject to their common unitary origin. Such an inequality is not obtained from Eqs.~\eqref{eq:budgetdmain} and \eqref{eq:fouriersharemain}; the all-dimensional projective optimality result of Ref.~\cite{TurtureanBao} is established by a different analytic method. The extension here isolates the geometric content and the additional constraint needed to recover that result by the present route.

\paragraph{Discussion.}
It is useful to separate three optimization statements that are often placed next to one another in the CGLMP literature. First, the unrestricted quantum maximum allows both the state and measurements to vary; for qutrits it is larger than Eq.~\eqref{eq:mainI}, is attained by a nonmaximally entangled state \cite{Acin2002}, and is now known exactly from sum-of-squares methods \cite{IoannouRosset}. Second, the problem solved here keeps the original CGLMP Bell functional and fixes the state to $\Phithree$; only the measurements vary. Third, one may instead redesign the Bell functional so that a maximally entangled state and the CGLMP measurement family jointly attain the global Tsirelson bound, as in the SATWAP construction \cite{SATWAP,Sarkar}. The present argument gives a geometric proof of the fixed-state statement for $d=3$, complementing the all-dimensional analytic result of Ref.~\cite{TurtureanBao}.

The proof also explains why the numerical optimum has been so stable across different searches. The four fractional phase offsets are not special because they belong to a convenient Fourier ansatz. They realize an equal decomposition of a fixed noncommutative obstruction. Any optimal strategy must reproduce the same angular share, the same commutator spectrum, and the same cyclic displacement profile on each of the four edges. This converts a numerical observation about one family of bases into a structural statement about the full qutrit projective measurement space.



\section*{Declaration of AI Assistance}

This work was carried out with substantial assistance from
ChatGPT (GPT-6). Under the author's direction, GPT-6 was
tasked with finding a rigorous proof of the conjectured
optimality of CGLMP measurements for maximally entangled
qutrits. GPT-6 contributed extensively to the development
of the proof, including the mathematical derivations,
the formulation and refinement of key arguments, and
consistency checks. Its assistance was used throughout
the preparation of the manuscript, from the initial
investigation to successive revisions. The author guided
the research, critically examined the arguments, and
reorganized and refined the proof to give it its present
form, exposition, and style. The author has independently
verified the mathematical results and assumes full
responsibility for their correctness and for the
contents of the manuscript.



\begin{thebibliography}{99}

\bibitem{CGLMP}
D. Collins, N. Gisin, N. Linden, S. Massar, and S. Popescu,
Bell inequalities for arbitrarily high-dimensional systems,
Phys. Rev. Lett. \textbf{88}, 040404 (2002).

\bibitem{Masanes}
L. Masanes,
Tight Bell inequality for $d$-outcome measurements correlations,
Quantum Inf. Comput. \textbf{3}, 345 (2003).

\bibitem{Durt2001}
T. Durt, D. Kaszlikowski, and M. \ifmmode\dot{Z}\else\.Z\fi{}ukowski,
Violations of local realism with quantum systems described by $N$-dimensional Hilbert spaces,
Phys. Rev. A \textbf{64}, 024101 (2001).

\bibitem{Acin2002}
A. Ac\'{i}n, T. Durt, N. Gisin, and J. I. Latorre,
Quantum nonlocality in two three-level systems,
Phys. Rev. A \textbf{65}, 052325 (2002).

\bibitem{AcinGillGisin}
A. Ac\'{i}n, R. Gill, and N. Gisin,
Optimal Bell tests do not require maximally entangled states,
Phys. Rev. Lett. \textbf{95}, 210402 (2005).

\bibitem{NPA}
M. Navascu\'{e}s, S. Pironio, and A. Ac\'{i}n,
Bounding the set of quantum correlations,
Phys. Rev. Lett. \textbf{98}, 010401 (2007).

\bibitem{ZG}
S. Zohren and R. D. Gill,
Maximal violation of the Collins-Gisin-Linden-Massar-Popescu inequality for infinite dimensional states,
Phys. Rev. Lett. \textbf{100}, 120406 (2008).

\bibitem{LVN}
B. Lang, T. V\'{e}rtesi, and M. Navascu\'{e}s,
Closed sets of correlations: Answers from the zoo,
J. Phys. A: Math. Theor. \textbf{47}, 424029 (2014).

\bibitem{SATWAP}
A. Salavrakos, R. Augusiak, J. Tura, P. Wittek, A. Ac\'{i}n, and S. Pironio,
Bell inequalities tailored to maximally entangled states,
Phys. Rev. Lett. \textbf{119}, 040402 (2017).

\bibitem{Sarkar}
S. Sarkar, D. Saha, J. Kaniewski, and R. Augusiak,
Self-testing quantum systems of arbitrary local dimension with minimal number of measurements,
npj Quantum Inf. \textbf{7}, 151 (2021).

\bibitem{IoannouRosset}
M. Ioannou and D. Rosset,
Noncommutative polynomial optimization under symmetry,
arXiv:2112.10803 (2021).

\bibitem{TurtureanBao}
D. Turturean and N. Bao,
The optimal CGLMP measurements on maximally entangled states,
unpublished preprint (2026), repository snapshot of 23 October 2026,
\href{https://github.com/physical-superintelligence/oqp-batch-2026/blob/7049357c42a8fe85262da08fdc10d57b4d9e67a0/oqp27-cglmp-inequalities/paper/oqp27-cglmp-inequalities.pdf}{GitHub repository snapshot}.

\bibitem{Thompson}
R. C. Thompson,
Proof of a conjectured exponential formula,
Linear Multilinear Algebra \textbf{19}, 187 (1986).

\bibitem{Antezana}
J. Antezana, G. Larotonda, and A. Varela,
Thompson-type formulae,
J. Funct. Anal. \textbf{262}, 1515 (2012).

\end{thebibliography}

\clearpage
\onecolumngrid
\begin{center}
{\large\bfseries Supplementary Material}
\end{center}

\setcounter{equation}{0}
\renewcommand{\theequation}{S\arabic{equation}}
\renewcommand{\theHequation}{S.\arabic{equation}}

\paragraph{Exact factorization of the measurement problem.}
Let $\mathsf A_x,\mathsf B_y\in U(d)$ collect the measurement basis vectors as columns and set $E_x=\overline{\mathsf A_x}$, $F_y=\mathsf B_y$. For the maximally entangled state,
\begin{equation}
p(a,b\vert x,y)=\frac1d\left|\left(E_x^\dagger F_y\right)_{ab}\right|^2.
\label{eqS:prob}
\end{equation}
Introduce
\begin{equation}
U_1=E_0^\dagger F_0,\quad
U_2=F_0^\dagger E_1,\quad
U_3=E_1^\dagger F_1,\quad
U_4=F_1^\dagger E_0S.
\label{eqS:U}
\end{equation}
The first three terms in Eq.~\eqref{eq:Bd} give $c_d(U_1)$, $c_d(U_2)$, and $c_d(U_3)$. For the last term, the $j$th column of $E_0S$ is the $(j-1)$th column of $E_0$, so $c_d(U_4)$ equals $\mathbb E[B_1-A_0-1]_d$. The product telescopes,
\begin{equation}
U_1U_2U_3U_4=S.
\label{eqS:telescope}
\end{equation}
Conversely, given any $U_1,U_2,U_3,U_4\in U(d)$ satisfying Eq.~\eqref{eqS:telescope}, choose
\begin{align}
E_0&=\id, & F_0&=U_1,\notag\\
E_1&=U_1U_2, & F_1&=U_1U_2U_3.
\end{align}
Then $F_1^\dagger E_0S=U_4$. Hence Eq.~\eqref{eq:factor} is an exact reformulation, not a relaxation.

\paragraph{Qutrit cost as an $SU(3)$ class function.}
For $U\in U(3)$ define the cyclic displacement weights
\begin{equation}
p_r(U)=\frac13\sum_{a-b=r\;\mathrm{mod}\;3}|U_{ab}|^2,
\qquad r=0,1,2.
\label{eqS:pr}
\end{equation}
They satisfy $p_0+p_1+p_2=1$ and $c_3(U)=p_1+2p_2$. With $K(U)$ from Eq.~\eqref{eq:K},
\begin{equation}
\frac13\Tr K(U)=p_0+\omega p_1+\omega^2p_2.
\label{eqS:traceK}
\end{equation}
Taking real and imaginary parts of Eq.~\eqref{eqS:traceK} and using normalization gives
\begin{equation}
\ell(K)=1-\operatorname{Re}\frac{\Tr K}{3}
-\frac1{\sqrt3}\operatorname{Im}\frac{\Tr K}{3},
\quad c_3(U)=\ell(K(U)).
\label{eqS:ell}
\end{equation}

The cocycle identity
\begin{equation}
K(UV)=K(U)\,U K(V)U^\dagger
\label{eqS:cocycle}
\end{equation}
is immediate from the definition. Let $P_j=U_1\cdots U_j$, $P_0=\id$, and
\begin{equation}
L_j=P_{j-1}K(U_j)P_{j-1}^\dagger.
\label{eqS:Lj}
\end{equation}
Repeated use of Eq.~\eqref{eqS:cocycle} gives
\begin{equation}
L_1L_2L_3L_4=K(U_1U_2U_3U_4)=K(S).
\end{equation}
Since $ZS=\omega^2SZ$, $K(S)=\omega^2\id$. Because $\ell$ depends only on the trace, it is invariant under conjugation and $c_3(U_j)=\ell(L_j)$. This proves Eq.~\eqref{eq:central}.

\paragraph{Angular gauge and subadditivity.}
Every $K\in SU(3)$ has a unique representative of its eigenangles in the Weyl alcove,
\begin{equation}
\theta_1\geq\theta_2\geq\theta_3,\qquad
\theta_1+\theta_2+\theta_3=0,\qquad
\theta_1-\theta_3\leq2\pi.
\label{eqS:alcove}
\end{equation}
For this representative, the quantity
\begin{equation}
m(K)=\min_{\substack{H=H^\dagger,\;\Tr H=0\\ \e^{iH}=K}}\lambda_{\max}(H)
\label{eqS:m}
\end{equation}
is $m(K)=\theta_1$. To see that no other traceless logarithm lowers the maximum eigenvalue, write any other logarithmic spectrum as $\theta_j+2\pi n_j$, with integers $n_j$ whose sum is zero. If every $n_j\leq0$, then all $n_j=0$. Otherwise some $n_j\geq1$, and the corresponding shifted angle is at least $\theta_3+2\pi\geq\theta_1$ by the alcove condition. Hence the canonical logarithm is optimal.

For Hermitian $A$ and $B$, Thompson's theorem gives unitaries $X,Y$ such that
\begin{equation}
\e^{iA}\e^{iB}=\e^{i(XAX^\dagger+YBY^\dagger)}.
\label{eqS:thompson}
\end{equation}
Taking optimal traceless logarithms of $K$ and $L$, the exponent on the right remains traceless. Therefore
\begin{align}
m(KL)
&\leq\lambda_{\max}(XAX^\dagger+YBY^\dagger)\notag\\
&\leq\lambda_{\max}(A)+\lambda_{\max}(B)
=m(K)+m(L).
\label{eqS:triangle}
\end{align}
This establishes the required subadditivity. The canonical angles of $\omega^2\id$ are $(4\pi/3,-2\pi/3,-2\pi/3)$, so $m(\omega^2\id)=4\pi/3$ and Eq.~\eqref{eq:angbudget} follows.

\paragraph{Sharp one-matrix lower envelope.}
Let
\begin{equation}
\phi(\theta)=1-\cos\theta-\frac1{\sqrt3}\sin\theta.
\label{eqS:phi}
\end{equation}
For the canonical angles in Eq.~\eqref{eqS:alcove},
\begin{equation}
\ell(K)=\frac13\sum_{j=1}^{3}\phi(\theta_j).
\label{eqS:ellangles}
\end{equation}
Fix $m=\theta_1$ and put $\theta_2=y$, $\theta_3=-m-y$. The alcove constraints become
\begin{equation}
-\frac m2\leq y\leq\min\{m,2\pi-2m\}.
\label{eqS:yint}
\end{equation}
The derivative of the $y$-dependent part is
\begin{equation}
\frac{d}{dy}\bigl[\phi(y)+\phi(-m-y)\bigr]
=\frac4{\sqrt3}\cos\!\left(\frac m2+\frac\pi6\right)
\sin\!\left(y+\frac m2\right).
\label{eqS:yder}
\end{equation}
For $0\leq m\leq2\pi/3$ the derivative is nonnegative on the allowed interval, so the minimum occurs at $y=-m/2$. For $2\pi/3\leq m\leq4\pi/3$ it is nonpositive, so the minimum occurs at $y=2\pi-2m$. Thus
\begin{equation}
\ell(K)\geq g(m(K)),
\label{eqS:ellg}
\end{equation}
where
\begin{equation}
g(m)=
\begin{cases}
\displaystyle g_1(m)=\frac{\phi(m)+2\phi(-m/2)}{3},
&0\leq m\leq\frac{2\pi}{3},\\[2mm]
\displaystyle g_2(m)=\frac{2\phi(m)+\phi(-2m)}{3},
&\frac{2\pi}{3}\leq m\leq\frac{4\pi}{3}.
\end{cases}
\label{eqS:g}
\end{equation}
The relevant derivatives are
\begin{align}
g_1'(m)&=\frac4{3\sqrt3}\cos\!\left(\frac m4-\frac\pi6\right)
\sin\frac{3m}{4},\notag\\
g_1''(m)&=\frac{\sqrt3}{9}\left[2\sin\!\left(m+\frac\pi3\right)
+\cos\!\left(\frac m2+\frac\pi6\right)\right],\notag\\
g_2'(m)&=\frac8{3\sqrt3}\cos\!\left(\frac m2+\frac\pi6\right)
\sin\frac{3m}{2}.
\label{eqS:gder}
\end{align}
Hence $g_1$ is increasing and convex on $[0,2\pi/3]$, while $g_2$ is increasing on $[2\pi/3,4\pi/3]$.

\paragraph{The scalar minimization.}
Write $m_j=m(L_j)$. Since $g$ is increasing, the minimum compatible with Eq.~\eqref{eq:angbudget} occurs at $\sum_jm_j=4\pi/3$. If all $m_j\leq2\pi/3$, convexity gives
\begin{equation}
\sum_{j=1}^{4}g_1(m_j)\geq4g_1(\pi/3).
\label{eqS:jensen}
\end{equation}
If one factor lies on the second branch, denote it by $M\geq2\pi/3$. There can be at most one such factor, and convexity of the other three gives
\begin{equation}
\sum_{j=1}^{4}g(m_j)\geq
H(M):=g_2(M)+3g_1\!\left(\frac{4\pi/3-M}{3}\right).
\label{eqS:H}
\end{equation}
It remains to show that $H(M)$ never beats the balanced point.

For $2\pi/3\leq M\leq5\pi/6$, monotonicity gives
\begin{equation}
H(M)\geq1+3g_1(\pi/6),
\end{equation}
and
\begin{equation}
1+3g_1(\pi/6)-4g_1(\pi/3)
=\frac43+\frac{2\sqrt3}{9}-\frac{2\sqrt6}{3}>0.
\label{eqS:r1}
\end{equation}
For $5\pi/6\leq M\leq\pi$, the argument of $g_1$ in Eq.~\eqref{eqS:H} lies in $[\pi/9,\pi/6]$. On $[0,\pi/6]$, Eq.~\eqref{eqS:gder} gives $g_1''\geq1/2$. Since $g_1(0)=g_1'(0)=0$, $g_1(x)\geq x^2/4$, and therefore
\begin{equation}
H(M)\geq g_2(5\pi/6)+\frac{\pi^2}{108}>4g_1(\pi/3).
\label{eqS:r2}
\end{equation}
Finally, for $\pi\leq M\leq4\pi/3$,
\begin{equation}
H(M)\geq g_2(\pi)=\frac43>4g_1(\pi/3).
\label{eqS:r3}
\end{equation}
Thus the balanced point is the global scalar minimum. Since
\begin{equation}
g_1(\pi/3)=\frac23\left(1-\frac1{\sqrt3}\right),
\end{equation}
Eq.~\eqref{eq:Bbound} follows.

\paragraph{Saturation by the standard phase settings.}
Let $F_{jk}=3^{-1/2}\omega^{jk}$ and $D(t)=\diag(1,\omega^t,\omega^{2t})$, so that $R(t)=FD(t)F^\dagger$ and $R(1)=S$. In a phase convention consistent with Eq.~\eqref{eq:Bd}, the standard CGLMP settings may be taken as
\begin{equation}
\alpha_0=0,\qquad \alpha_1=-\frac12,\qquad
\beta_0=\frac14,\qquad \beta_1=\frac34.
\label{eqS:phases}
\end{equation}
Substitution in Eq.~\eqref{eqS:U} gives $U_j=R(1/4)$ for all four steps. Moreover,
\begin{equation}
\spec K(R(1/4))=
\left\{\e^{i\pi/3},\e^{-i\pi/6},\e^{-i\pi/6}\right\}.
\label{eqS:spectrum}
\end{equation}
Thus $m=\pi/3$ for each step, the first branch of Eq.~\eqref{eqS:g} is saturated, and
\begin{equation}
B_3=4g_1(\pi/3)=\frac83\left(1-\frac1{\sqrt3}\right).
\end{equation}


\paragraph{General-dimensional harmonic form.}
For completeness, define
\begin{equation}
Z_d=\diag(1,\omega_d,\ldots,\omega_d^{d-1}),
\qquad \omega_d=\e^{2\pi i/d},
\end{equation}
and
\begin{equation}
t_k(U)=\frac1d\Tr\!\left(Z_d^kUZ_d^{-k}U^\dagger\right).
\end{equation}
If $p_r=d^{-1}\sum_{a-b=r}|U_{ab}|^2$, then $t_k=\sum_rp_r\omega_d^{kr}$. Fourier inversion and
\begin{equation}
\sum_{r=0}^{d-1}r\,q^r=-\frac{d}{1-q},\qquad q^d=1,\quad q\neq1,
\end{equation}
give the exact identity
\begin{equation}
c_d(U)=\frac{d-1}{2}
-\sum_{k=1}^{d-1}\frac{t_k(U)}{1-\omega_d^{-k}}.
\label{eqS:general}
\end{equation}
The terms $k$ and $d-k$ are conjugate, so the right-hand side is real. Each matrix
\begin{equation}
K_k(U)=Z_d^kUZ_d^{-k}U^\dagger
\end{equation}
obeys the same cocycle law as Eq.~\eqref{eqS:cocycle}, and the endpoint imposed by the shift is $K_k(S)=\omega_d^{-k}\id$. For $d=3$, $t_2=\overline{t_1}$ and the cost is controlled by one complex harmonic. For $d\geq4$, several independent harmonics remain, which is the obstruction to repeating the scalar $SU(3)$ argument directly.



\paragraph{Simultaneous angular budgets for all dimensions.}
The harmonic commutators defined above have a common multiplicative origin. For $1\leq k\leq d-1$, let $P_0=\id$, $P_j=U_1\cdots U_j$, and put
\begin{equation}
L_{j,k}=P_{j-1}K_k(U_j)P_{j-1}^{\dagger}.
\label{eqS:Ljk}
\end{equation}
The cocycle identity gives the exact central-product constraints
\begin{equation}
\prod_{j=1}^{4}L_{j,k}=K_k(S)=\omega_d^{-k}\id.
\label{eqS:centralall}
\end{equation}
We use the same angular gauge $m$ as in Eq.~\eqref{eqS:m}, now on $SU(d)$: it is the minimum largest eigenvalue among traceless Hermitian logarithms. A traceless logarithm whose ordered eigenangles have spread at most $2\pi$ attains this minimum. Indeed, if another logarithm lowers its maximum, it must shift some eigenangle upwards by $2\pi$ while shifting others down; the spread condition rules out a lower maximum. Thompson's theorem again yields
$m(AB)\leq m(A)+m(B)$.
For $\omega_d^{-k}\id$, a minimizing traceless logarithm has $k$ angles $2\pi(1-k/d)$ and $d-k$ angles $-2\pi k/d$. Therefore every admissible unitary loop obeys the \emph{simultaneous angular budgets}
\begin{equation}
\sum_{j=1}^{4}m\bigl(K_k(U_j)\bigr)\geq
2\pi\left(1-\frac{k}{d}\right),\quad k=1,\ldots,d-1.
\label{eqS:allbudgets}
\end{equation}
At $d=3$ and $k=1$, this recovers Eq.~\eqref{eq:angbudget}. For example, in dimension four, the right sides for $k=1,2,3$ are $3\pi/2$, $\pi$, and $\pi/2$ respectively. These constraints hold for arbitrary projective bases and require no Fourier ansatz.

\paragraph{The Fourier path saturates all angular budgets.}
Let $F_{ab}=d^{-1/2}\omega_d^{ab}$ and define
\begin{equation}
R_d(t)=F\diag(1,\omega_d^t,\ldots,\omega_d^{(d-1)t})F^\dagger.
\label{eqS:Rgeneral}
\end{equation}
Then $R_d(s)R_d(t)=R_d(s+t)$ and $R_d(1)=S$. In the Fourier basis, $Z_d$ becomes a cyclic permutation, so $K_k(R_d(t))$ is diagonal with eigenvalues
\begin{equation}
\spec K_k(R_d(t))=\left\{
\begin{aligned}
&e^{2\pi it(1-k/d)}&&\text{($k$ times)},\\
&e^{-2\pi itk/d}&&\text{($d-k$ times)}
\end{aligned}\right\}.
\label{eqS:specall}
\end{equation}
For $0\leq t\leq1$, these eigenangles have vanishing sum and spread $2\pi t\leq2\pi$, hence
\begin{equation}
m\bigl(K_k(R_d(t))\bigr)=2\pi t\left(1-\frac{k}{d}\right).
\label{eqS:mR}
\end{equation}
The equal-step choice $U_j=R_d(1/4)$ consequently saturates \emph{every} inequality in Eq.~\eqref{eqS:allbudgets}. So does any positive subdivision $U_j=R_d(t_j)$ with $t_j\geq0$ and $\sum_j t_j=1$. Thus simultaneous saturation of the budgets alone cannot distinguish equal steps from unequal ones.

\paragraph{Equal sharing along the Fourier path.}
The cost restricted to Eq.~\eqref{eqS:Rgeneral} admits the exact trigonometric expression
\begin{equation}
\begin{split}
f_d(t):=c_d(R_d(t))=\frac{d-1}{2}-\frac1d\sum_{r=1}^{d-1}
(d-r)\frac{\sin[(1-2t)\pi r/d]}{\sin(\pi r/d)}.
\end{split}
\label{eqS:ft}
\end{equation}
It follows by evaluating the geometric series for the matrix entries of $R_d(t)$ or by substituting into Eq.~\eqref{eqS:general}. Differentiation gives
\begin{equation}
f_d''(t)=\frac{4\pi^2}{d^3}\sum_{r=1}^{d-1}r^2(d-r)
\frac{\sin[(1-2t)\pi r/d]}{\sin(\pi r/d)}.
\label{eqS:fsecond}
\end{equation}
Consequently, $f_d$ is strictly convex for $0<t<1/2$ but strictly concave for $1/2<t<1$. Global equal sharing cannot be concluded from convexity alone.

The required inequality is elementary but stronger: if $N\geq3$, $0<\varphi<\pi$, $t_j\geq0$ and $\sum_{j=1}^{N}t_j=1$, then
\begin{equation}
\sum_{j=1}^{N}\sin[(1-2t_j)\varphi]\leq
N\sin[(1-2/N)\varphi].
\label{eqS:sineav}
\end{equation}
To prove it, set $x_j=1-2t_j$, so $\sum_jx_j=N-2$. If all $x_j\geq0$, concavity of $\sin(\varphi x)$ on $[0,1]$ gives the claim. Otherwise exactly one is negative; write it $x_1=-a$, $0<a\leq1$. Concavity for the other $N-1$ terms gives the upper bound
\begin{equation}
-\sin(a\varphi)+(N-1)\sin\!\left(\frac{N-2+a}{N-1}\varphi\right).
\label{eqS:h}
\end{equation}
This bound decreases with $a$, since its derivative is
\begin{equation}
\varphi\left[\cos\!\left(\frac{N-2+a}{N-1}\varphi\right)-\cos(a\varphi)\right]\leq0.
\end{equation}
At $a=0$ the previous concavity argument applies. Equality for $N\geq3$ and $0<\varphi<\pi$ requires $t_1=\cdots=t_N=1/N$.

Applying Eq.~\eqref{eqS:sineav} term by term in Eq.~\eqref{eqS:ft}, whose sine coefficients have positive weights, establishes the all-dimensional \emph{balanced Fourier-subdivision theorem}
\begin{equation}
\sum_{j=1}^{N}c_d(R_d(t_j))\geq
N c_d\bigl(R_d(1/N)\bigr),\quad
\sum_jt_j=1,\ t_j\geq0,\ N\geq3.
\label{eqS:balancedN}
\end{equation}
In particular, for the four-setting CGLMP cycle,
\begin{equation}
\begin{split}
B_d^{\mathrm F}&=4c_d(R_d(1/4))\\
&=2(d-1)-\frac2d\sum_{r=1}^{d-1}(d-r)
\sec\!\left(\frac{\pi r}{2d}\right).
\end{split}
\label{eqS:BF}
\end{equation}
Equation~\eqref{eqS:balancedN} proves the optimality of the equal-step factorization \emph{within the positive Fourier path}, not among arbitrary unitary loops.

\paragraph{Why the higher-dimensional cost cannot depend on one commutator spectrum.}
Already in $d=4$, two unitaries with the same $K_1$ spectrum can have different costs. Write
$X_{ab}=\ket a\bra b+\ket b\bra a$ and take
\begin{equation}
U=e^{i(\pi/4)X_{01}},\qquad V=e^{i(\pi/6)X_{02}}.
\end{equation}
For a two-level rotation $e^{i\theta X_{ab}}$, the nontrivial eigenvalues of $K_1$ are $e^{\pm i\chi}$, where
$\sin(\chi/2)=|\sin\theta\sin[\pi(a-b)/4]|$.
Thus
\begin{equation}
\begin{gathered}
\spec K_1(U)=\spec K_1(V)
=\{e^{i\pi/3},e^{-i\pi/3},1,1\},\\
c_4(U)=\frac12,\qquad c_4(V)=\frac14.
\end{gathered}
\label{eqS:d4ex}
\end{equation}
The second harmonic separates them: $K_2(V)=\id$ while $K_2(U)$ has nontrivial eigenvalues $e^{\pm i\pi/2}$. This rules out a direct repetition of the qutrit spectral-envelope argument using $K_1$ alone.

Nor can the harmonic contributions be optimized independently. In $d=4$,
\begin{equation}
c_4(U)=\frac32-\operatorname{Re}t_1(U)
-\operatorname{Im}t_1(U)-\frac12t_2(U),
\label{eqS:d4harm}
\end{equation}
with $t_2(U)$ real. For instance, $U_j=R_4(13/4)$ for all four steps satisfies $\prod_jU_j=R_4(13)=S$ and improves the first-harmonic contribution relative to $R_4(1/4)$, while worsening the second by a larger amount. The harmonics must be controlled jointly.

\paragraph{The remaining coupled inequality.}
The higher harmonics are not independent matrices. Direct multiplication gives
\begin{equation}
K_k(U)=\bigl(Z_d^{k-1}K_1(U)Z_d^{-(k-1)}\bigr)\cdots
\bigl(Z_dK_1(U)Z_d^{-1}\bigr)K_1(U).
\label{eqS:Kiterate}
\end{equation}
The exact global cost inequality needed to complete a geometric proof for any $d$ is
\begin{equation}
\begin{split}
&\sum_{j=1}^{4}\sum_{k=1}^{d-1}
\operatorname{Re}\!\left(\frac{t_k(U_j)}{1-\omega_d^{-k}}\right)\\
&\qquad\leq4\sum_{k=1}^{d-1}
\operatorname{Re}\!\left(\frac{t_k(R_d(1/4))}{1-\omega_d^{-k}}\right),
\quad \prod_jU_j=S.
\end{split}
\label{eqS:missing}
\end{equation}
This is equivalent to $B_d\geq B_d^{\mathrm F}$ and hence is \emph{not} a separately proved lemma. The all-dimensional optimality theorem of Ref.~\cite{TurtureanBao} establishes its validity by a different analytic route; the angular budgets and Fourier-path theorem above do not constitute an independent proof of this inequality. An alternative geometric formulation uses the conjugated clock observables $Q_j=P_jZ_dP_j^\dagger$. They satisfy $Q_0=Z_d$, $Q_4=\omega_d Z_d$, and
\begin{equation}
t_k(U_j)=\frac1d\Tr(Q_{j-1}^{k}Q_j^{-k}).
\label{eqS:orbit}
\end{equation}
The missing ingredient is therefore a sharp, joint four-edge cost inequality on this twisted unitary orbit, rather than a separate estimate for each Fourier harmonic.

\end{document}
